%% kapa: exam=1 round=31 year=2020 period=1 form=A subject=real_estate subject_ko=부동산학원론 seq=32 number=112
%% source: Q-Net 감정평가사 기출문제 (cst00309) artlSeq=5207219 file=extracted/1008/5207219/A형/A형/2020년도 제31회 감정평가사 1차 시험_1교시_A형.pdf
%% conversion: kordoc 4.24.51 -> tools/gen1.py (auto); review=proofread-sonnet; flags=none
\begin{question}{112}
토지에 관한 설명으로 옳지 않은 것은?
\begin{choices}
  \item 빈지는 일반적으로 바다와 육지 사이의 해변 토지와 같이 소유권이 인정되며 이용실익이 있는 토지이다.
  \item 맹지는 타인의 토지에 둘러싸여 도로에 어떤 접속면도 가지지 못하는 토지이며, 건축법에 의해 원칙적으로 건물을 세울 수 없다.
  \item 법지는 택지경계와 접한 경사된 토지부분과 같이 법률상으로는 소유를 하고 있지만 이용실익이 없는 토지이다.
  \item 후보지는 부동산의 주된 용도적 지역인 택지지역, 농지지역, 임지지역 상호간에 전환되고 있는 지역의 토지이다.
  \item 이행지는 부동산의 주된 용도적 지역인 택지지역, 농지지역, 임지지역의 세분된 지역 내에서 용도전환이 이루어지고 있는 토지이다.
\end{choices}
\end{question}
\answer{1}
